\section{Robustness of the Joint Job Loss}\label{app:robust}

This appendix documents the checks summarised in Section~\ref{sec:jointloss}: the CPS monthly checks on reporting,
temporary layoffs, the season, persistence, class of worker, a common employer and attrition; the replication in the
Displaced Worker Supplement; and the replication and the added-worker path in the SIPP.

\subsection{The CPS monthly checks}

Three features of the CPS could produce a joint job loss that is not one: one household member usually reports for
both spouses, half of the husbands' job losses are temporary layoffs, and some jobs end with the season or the school
year. None of them accounts for the fact; Table~\ref{tab:cps_checks} reports the estimates.\footnote{Linear probability models with state $\times$ month fixed effects,
standard errors clustered by household; the respondent identifiers (PUSLFPRX, HURESPLI) are matched for 99.6\% of
couple-months. Code: \texttt{IGE/37.CPSJointLossChecks.py} (\texttt{output/cps\_joint\_loss\_checks.md}).}

\emph{Who reports.} The excess is 3.8 points when the wife reports for herself and the husband's job loss is reported by
proxy, and 3.8 points in the reverse case; the difference is 0.0003 ($p=0.936$), and it is the same when the same person
reports in both months. Couples in which a third household member reports for both spouses (4\% of couples) show a
larger excess (0.093$^{***}$); excluding them leaves 3.8 points. Couples in which both spouses report for themselves
cannot be formed, because the CPS takes one respondent per household and month.

\emph{Temporary layoffs, the season and persistence.} Excluding the husband's temporary layoffs, the excess is 2.7
points (0.027$^{***}$); excluding the wife's as well, 2.6 points, on a base of 0.23\%, so the wife's risk is 14 times the
base. A temporary layoff of the husband followed by a permanent job loss of the wife adds only 0.4 points. Temporary
layoffs therefore account for about a third of the headline excess. Excluding June to September leaves 3.7 points, and
excluding the pandemic years 3.5 points. When both spouses are still unemployed two months later, the excess is 2.5
points (0.025$^{***}$) on a base of 0.20\%. The strictest version---both spouses still unemployed two months later, the
husband's loss not a temporary layoff, outside June to September and outside 2020--21---gives 1.6 points
(0.016$^{***}$) on a base of 0.17\%: a joint permanent job loss is about 13 times as likely in the month the husband loses
his job as in other months. The excess is positive in all ten major industries. In education and health it is large
overall (0.091$^{***}$) but 2.0 points for the husband's permanent losses, consistent with school-year layoffs; couples
who both work for government or are both self-employed show the largest excess (0.153$^{***}$ and 0.159$^{***}$), the
latter likely family businesses, and couples who both work in the private sector show 3.0 points for the husband's
permanent losses.

\emph{A common employer.} The public CPS has no employer identifier. The tightest proxy, couples in the same detailed
industry, the same detailed occupation and the same class of worker (6.3\% of couples), shows a joint-loss probability
of 35.7\% (excess 0.337$^{***}$) and accounts for 35\% of the excess; all other couples show 2.6 points. If 10--12\% of
dual-earner couples work at the same establishment \citep{hyatt2019coworking}, the excess requires a joint-loss
probability of 21--41\% among them, depending on the excess elsewhere. A shared employer plausibly carries a third to a
half of the excess. Couples observed in all four months of the rotation show 3.8 points, so attrition after a job loss
does not create the excess.

\begin{table}[!htbp]
\centering
\caption{Robustness of the same-month joint job loss (CPS basic monthly, 1994--2023)}
\label{tab:cps_checks}
\resizebox{\ifdim\width>\textwidth\textwidth\else\width\fi}{!}{%
\resizebox*{!}{\ifdim\totalheight>0.78\textheight 0.78\textheight\else\totalheight\fi}{%
\begin{tabular}{lccccc}
\toprule
Sample & \begin{tabular}[b]{@{}c@{}}Wife loses\\her job\end{tabular} & \begin{tabular}[b]{@{}c@{}}Wife loses her job,\\excluding 2020--21\end{tabular} & \begin{tabular}[b]{@{}c@{}}Wife's loss not a\\temporary layoff\end{tabular} & \begin{tabular}[b]{@{}c@{}}Wife still\\unemployed at $t+2$\end{tabular} & \begin{tabular}[b]{@{}c@{}}Base rate,\\\% a month\end{tabular} \\
 & (1) & (2) & (3) & (4) & (5) \\
\midrule
\multicolumn{6}{l}{\textit{A. Who reports (month of the husband's job loss)}} \\
All couples & $0.040^{***}$ & $0.035^{***}$ &  &  & 0.54 \\
 & (0.002) & (0.001) &  &  & \\
Wife reports for herself, husband by proxy & $0.038^{***}$ & $0.034^{***}$ &  &  & 0.60 \\
 & (0.002) & (0.002) &  &  & \\
\quad and the wife reports his job loss & $0.039^{***}$ & $0.034^{***}$ &  &  & 0.60 \\
 & (0.002) & (0.002) &  &  & \\
Husband reports for himself, wife by proxy & $0.038^{***}$ & $0.032^{***}$ &  &  & 0.47 \\
 & (0.002) & (0.002) &  &  & \\
\quad and the husband reports her job loss & $0.039^{***}$ & $0.032^{***}$ &  &  & 0.47 \\
 & (0.002) & (0.002) &  &  & \\
Wife self, husband proxy in both months & $0.037^{***}$ & $0.031^{***}$ &  &  & 0.59 \\
 & (0.002) & (0.002) &  &  & \\
Husband self, wife proxy in both months & $0.037^{***}$ & $0.030^{***}$ &  &  & 0.46 \\
 & (0.003) & (0.002) &  &  & \\
A third member reports both & $0.093^{***}$ & $0.084^{***}$ &  &  & 0.54 \\
 & (0.010) & (0.010) &  &  & \\
Third-member reports excluded & $0.038^{***}$ & $0.033^{***}$ &  &  & 0.54 \\
 & (0.002) & (0.001) &  &  & \\
\addlinespace
\multicolumn{6}{l}{\textit{B. Temporary layoffs and the season}} \\
All husband job losses & $0.040^{***}$ & $0.035^{***}$ & $0.015^{***}$ & $0.015^{***}$ & 0.54 \\
 & (0.002) & (0.001) & (0.001) & (0.001) & \\
His loss not a temporary layoff & $0.027^{***}$ & $0.027^{***}$ & $0.026^{***}$ & $0.012^{***}$ & 0.54 \\
 & (0.002) & (0.002) & (0.002) & (0.001) & \\
His loss a temporary layoff & $0.054^{***}$ & $0.044^{***}$ & $0.004^{***}$ & $0.017^{***}$ & 0.54 \\
 & (0.002) & (0.002) & (0.001) & (0.002) & \\
Not June--September & $0.037^{***}$ & $0.031^{***}$ & $0.015^{***}$ & $0.015^{***}$ & 0.48 \\
 & (0.002) & (0.002) & (0.001) & (0.001) & \\
Not a temporary layoff, not June--September & $0.026^{***}$ & $0.025^{***}$ & $0.025^{***}$ & $0.012^{***}$ & 0.48 \\
 & (0.002) & (0.002) & (0.002) & (0.002) & \\
\addlinespace
\multicolumn{6}{l}{\textit{C. Persistent job losses (both observed two months later)}} \\
All husband job losses & $0.038^{***}$ & $0.034^{***}$ & $0.014^{***}$ & $0.015^{***}$ & 0.54 \\
 & (0.002) & (0.002) & (0.001) & (0.001) & \\
His loss persistent & $0.033^{***}$ & $0.028^{***}$ & $0.017^{***}$ & $0.025^{***}$ & 0.54 \\
 & (0.003) & (0.002) & (0.002) & (0.002) & \\
Persistent, not a temporary layoff & $0.027^{***}$ & $0.026^{***}$ & $0.025^{***}$ & $0.018^{***}$ & 0.54 \\
 & (0.003) & (0.003) & (0.003) & (0.002) & \\
Persistent, not a temporary layoff, not June--September & $0.028^{***}$ & $0.026^{***}$ & $0.026^{***}$ & $0.017^{***}$ & 0.47 \\
 & (0.004) & (0.004) & (0.003) & (0.003) & \\
\addlinespace
\multicolumn{6}{l}{\textit{D. Class of worker}} \\
Both private & $0.041^{***}$ & $0.036^{***}$ & $0.017^{***}$ & $0.017^{***}$ & 0.59 \\
 & (0.002) & (0.002) & (0.001) & (0.002) & \\
Both private, his loss not a temporary layoff & $0.030^{***}$ & $0.030^{***}$ & $0.029^{***}$ & $0.015^{***}$ & 0.59 \\
 & (0.002) & (0.002) & (0.002) & (0.002) & \\
Both government & $0.153^{***}$ & $0.156^{***}$ & $0.013^{**}$ & $0.040^{***}$ & 0.38 \\
 & (0.016) & (0.017) & (0.006) & (0.010) & \\
Both self-employed & $0.159^{***}$ & $0.142^{***}$ & $0.073^{***}$ & $0.043^{***}$ & 0.30 \\
 & (0.016) & (0.017) & (0.013) & (0.012) & \\
Husband private, wife government & $0.013^{***}$ & $0.007^{***}$ & $0.002^{*}$ & $0.004^{**}$ & 0.60 \\
 & (0.003) & (0.002) & (0.001) & (0.002) & \\
Husband government, wife private & $0.027^{***}$ & $0.024^{***}$ & $0.008^{**}$ & $0.006$ & 0.46 \\
 & (0.007) & (0.006) & (0.004) & (0.004) & \\
\addlinespace
\multicolumn{6}{l}{\textit{E. Shared workplace}} \\
Same major industry & $0.166^{***}$ &  &  &  & 0.44 \\
 & (0.007) &  &  &  & \\
Same detailed industry & $0.211^{***}$ &  &  &  & 0.41 \\
 & (0.009) &  &  &  & \\
Same detailed industry and occupation & $0.315^{***}$ &  &  &  & 0.41 \\
 & (0.014) &  &  &  & \\
Same detailed industry, occupation and class & $0.337^{***}$ &  &  &  & 0.41 \\
 & (0.015) &  &  &  & \\
All other couples & $0.026^{***}$ &  &  &  & 0.55 \\
 & (0.001) &  &  &  & \\
\addlinespace
\multicolumn{6}{l}{\textit{F. Attrition}} \\
Couples observed in all four months & $0.038^{***}$ & $0.033^{***}$ &  &  & 0.54 \\
 & (0.002) & (0.002) &  &  & \\
Couples not observed in all four months & $0.052^{***}$ & $0.045^{***}$ &  &  & 0.57 \\
 & (0.004) & (0.004) &  &  & \\
\bottomrule
\end{tabular}}}
\par\smallskip
\begin{minipage}{\textwidth}\footnotesize \textit{Notes:} Each cell is the coefficient on the husband losing his job in a linear probability model of the outcome in the column header, among married couples in which both spouses work in the previous month, with state $\times$ month fixed effects; standard errors clustered by household (state clustering gives the same significance). Column (5): the wife's job-loss rate when the husband does not lose his job. A job loser is unemployed after a layoff, a job loss or the end of a temporary job, having been employed the month before. Persistent: still unemployed two months later. Panel A: labor force information collected from the person (self) or from another household member (proxy); couples in which both spouses report for themselves cannot be formed because the CPS takes one respondent per household and month. $^{***}$ $p<0.01$, $^{**}$ $p<0.05$, $^{*}$ $p<0.1$. Source: \texttt{IGE/37.CPSJointLossChecks.py} (\texttt{output/cps\_joint\_loss\_checks.csv}); table built by \texttt{IGE/48.PaperATables.py}.\end{minipage}
\end{table}


\subsection{The Displaced Worker Supplement}

The CPS Displaced Worker Supplement, fielded every two years from 1984 to 2024, asks every adult whether he or she lost
a job in the last three years because the plant closed or moved, there was insufficient work or the position was
abolished, and in which year. It combines the reason-based measure of the PSID with the sample size of the CPS: 417,269
couples in which both spouses answer the supplement, 24,038 displaced husbands and 1,968 couples in which both were
displaced.\footnote{IPUMS CPS, 21 supplements; couples aged 20--64 naming each other as spouse; the 1984--1992 five-year
window and reasons are harmonised to the later three-year definition. Linear probability models with state $\times$
survey-year fixed effects and controls for both spouses' age, education and race; standard errors clustered by
household (state clustering gives the same significance). Code: \texttt{IGE/39.CPSDisplacedWorkers.py}
(\texttt{output/cps\_dws\_joint\_loss.md}).} A wife was displaced in the three-year window in 8.2\% of couples in which
the husband was displaced and in 3.3\% of the others (an excess of 4.4 points, 0.044$^{***}$, s.e.\ 0.002), and in the same
calendar year in 4.7\% against 1.1\% (0.033$^{***}$, s.e.\ 0.001), while a displacement of the husband in another year of
the window adds only 0.5 points. Both spouses lost their jobs to plant closings 4.7 times as often as independent
closings would imply (0.046$^{***}$), and 9.8 times within the same year. Joint displacements share the lost job's
detailed industry in 29\% of cases, 41\% within the same year and 68\% for same-year joint closings, against 1.3\% for
randomly paired spouses; the excess rises by 40 points when the spouses worked in the same detailed industry
(0.399$^{***}$). The tetrachoric correlation of the spouses' displacements is 0.26 within the same year (0.35 for plant
closings), close to the PSID's 0.24. The excess is the same for Black and White couples ($-0.001$, $p=0.894$), in
1984--1998 and 2000--2024 (difference 0.004, $p=0.252$) and in recession and expansion windows (0.004, $p=0.305$). After
a joint displacement neither spouse works at the survey more often: 8.8 points more (0.088$^{***}$) on a base of 9.3\%.

\subsection{The SIPP}

The public-use SIPP core files of the 2004 (waves 1--12, 2003--2007), 2008 (waves 1--16, 2008--2013) and 2014
(waves 1--4, 2013--2016) panels give 1,595,445 couple-months of 50,692 married couples aged 20--60 who name each other
as spouse. An involuntary separation is an employer job that ends for layoff or slack work, closure, discharge or the end
of a temporary job, from the reason-for-leaving item. The public employer numbers are person-specific, so a shared
employer cannot be identified.\footnote{Linear probability models among couples in which both spouses were employed the
month before, with state $\times$ panel $\times$ rotation $\times$ month fixed effects; standard errors clustered by
household. Code: \texttt{IGE/40.SIPPJointLoss.py} (\texttt{output/sipp\_joint\_loss.md}).}

The wife's involuntary separation rate is 3.28\% in the month in which her husband separates and 0.43\% otherwise
(ratio 7.7; excess 0.027$^{***}$, s.e.\ 0.003, on 4,479 husbands' separations and 143 joint events); requiring that she has
no job the next month, as in the CPS, gives a ratio of 13.1 (excess 0.028$^{***}$). Over twelve months the ratio is 2.4,
against 2.6 in the PSID's couple-years. The excess holds outside the interview seam (0.027$^{***}$), for couples in which
both spouses report for themselves (0.019$^{***}$), without imputed records (0.028$^{***}$), in each panel (2004
0.036$^{***}$, 2008 0.022$^{***}$, 2014 0.028$^{***}$) and outside the 2008--09 recession (0.030$^{***}$). It rises by 22
points when the spouses share a detailed industry (0.220$^{***}$), where 23.9\% of wives separate in their husband's
month, and by 11 points in the same major industry. The timing matches the CPS: the excess is 2.6 points in the month of
the husband's separation, 0.4 points in the month before, and spread thinly over the months around it.

The SIPP differs from the CPS and the Displaced Worker Supplement in one respect: for Black husbands the same-month
excess is 0.007 ($p=0.123$) against 0.027$^{***}$ for White husbands (difference $-0.017^{***}$), on six joint events among
435 Black husbands' separations; over the month of the separation and the following twelve the excess is similar
(Black 0.050$^{***}$, White 0.063$^{***}$). With so few events this difference is fragile.

The SIPP also traces the wife's response month by month for three years. Among wives with no job the month before, the
probability of looking for work rises at once (0.053$^{***}$) and fades over three years (0.013, $p=0.246$ at 36 months),
while employment builds up slowly: 0.019$^{***}$ in the month of the husband's separation, 0.037$^{***}$ a month later and
0.054$^{***}$ after six months, and still 0.044$^{***}$ after 24 months. These wives were more often employed and looking
six months earlier, so some had recently lost or left jobs; among wives with no job in any of the six previous months,
employment rises from 0.015$^{***}$ to 0.045$^{***}$ after six months. Wives who were working see their employment fall by
3.4 points a month after the husband's separation, about half of it their own job loss, so the net change in all wives'
employment is about zero for a year ($-0.007^{**}$ after one month, 0.006, $p=0.297$, after twelve); leaving out the wives
who lost their own job it is 0.011$^{*}$ after twelve months and 0.020$^{**}$ after 24.
